\documentclass[border=9pt]{standalone}
\input{figure_preamble.tex}

\begin{document}

\begin{tikzpicture}

% figure_style.tex 会在 background 层绘制整张图的浅色底板。
% 增加 highlight 层，使右上角的边框和色块位于底板之上、文字之下。
\pgfdeclarelayer{highlight}
\pgfsetlayers{background,highlight,main}

\def\sep{0.6}
\def\stepsize{0.4cm}
\def\ycoordp{4}
\def\vectorsize{1cm}
\def\smallsize{0.13}

\tikzstyle{snode}=[
  fig/module,
  minimum width=2.5cm,
  minimum height=0.6cm
]

\tikzstyle{cnode}=[
  minimum width=0.3cm,
  minimum height=\vectorsize
]

\begin{scope}

% =========================================================
% 当前序列 x
% =========================================================

\node[anchor=center] (x1) at (0,0)
  {\texttt{A}};

\node[anchor=center] (x2) at ([xshift=\stepsize]x1)
  {\texttt{D}};

\node[anchor=center] (x3) at ([xshift=\stepsize]x2)
  {\texttt{B}};

\node[anchor=center] (x4) at ([xshift=\stepsize]x3)
  {\texttt{C}};

\node[anchor=center] (x5) at ([xshift=\stepsize]x4)
  {\texttt{C}};


% =========================================================
% 下方中间序列 y_2
% =========================================================

\node[anchor=center] (y21) at ([yshift=-2cm]x1)
  {\texttt{B}};

\node[anchor=center] (y22) at ([xshift=\stepsize]y21)
  {\texttt{D}};

\node[anchor=center] (y23) at ([xshift=\stepsize]y22)
  {\texttt{B}};

\node[anchor=center] (y24) at ([xshift=\stepsize]y23)
  {\texttt{C}};

\node[anchor=center] (y25) at ([xshift=\stepsize]y24)
  {\texttt{C}};


% =========================================================
% 下方右侧序列 y_3
% =========================================================

\node[anchor=west] (y31) at ([xshift=0.8cm]y25.east)
  {\texttt{A}};

\node[anchor=center] (y32) at ([xshift=\stepsize]y31)
  {\texttt{D}};

\node[anchor=center] (y33) at ([xshift=\stepsize]y32)
  {\texttt{B}};

\node[anchor=center] (y34) at ([xshift=\stepsize]y33)
  {\texttt{D}};

\node[anchor=center] (y35) at ([xshift=\stepsize]y34)
  {\texttt{C}};


% =========================================================
% 下方左侧序列 y_1
% =========================================================

\node[anchor=east] (y15) at ([xshift=-0.8cm]y21.west)
  {\texttt{C}};

\node[anchor=center] (y14) at ([xshift=-\stepsize]y15)
  {\texttt{C}};

\node[anchor=center] (y13) at ([xshift=-\stepsize]y14)
  {\texttt{B}};

\node[anchor=center] (y12) at ([xshift=-\stepsize]y13)
  {\texttt{C}};

\node[anchor=center] (y11) at ([xshift=-\stepsize]y12)
  {\texttt{A}};


% =========================================================
% 序列外框和标签
% =========================================================

\begin{pgfonlayer}{background}

\node[snode] (xbox) at (x3) {};

\node[snode] (ybox1) at (y13) {};

\node[snode] (ybox2) at (y23) {};

\node[snode] (ybox3) at (y33) {};

\node[anchor=east] (xlabel)
  at ([xshift=-0.2cm]xbox.west)
  {$\mathbf{x}$};

\node[anchor=north] (y1label)
  at ([yshift=-0.1cm]ybox1.south)
  {$\mathbf{y}_1$};

\node[anchor=north] (y2label)
  at ([yshift=-0.1cm]ybox2.south)
  {$\mathbf{y}_2$};

\node[anchor=north] (y3label)
  at ([yshift=-0.1cm]ybox3.south)
  {$\mathbf{y}_3$};

\end{pgfonlayer}


% 两侧省略号

\node[anchor=east] (y0)
  at ([xshift=-0.2cm]ybox1.west)
  {$\cdots$};

\node[anchor=west] (y4)
  at ([xshift=0.2cm]ybox3.east)
  {$\cdots$};


% =========================================================
% 下方序列指向当前序列
% =========================================================

\draw[->]
  ([yshift=1pt]ybox1.north)
  --
  ([xshift=-0.5cm,yshift=-2pt]xbox.south);

\draw[->]
  ([yshift=1pt]ybox2.north)
  --
  ([xshift=0cm,yshift=-1pt]xbox.south);

\draw[->]
  ([yshift=1pt]ybox3.north)
  --
  ([xshift=0.5cm,yshift=-2pt]xbox.south);


% =========================================================
% 左侧时间轴
% =========================================================

\draw[->,thick]
  ([xshift=-4.5cm,yshift=-3.0cm]x1.west)
  --
  ([xshift=-4.5cm,yshift=1.4cm]x1.west);

\node[anchor=east] (timelabel)
  at ([xshift=-4.6cm,yshift=1.2cm]x1.west)
  {\footnotesize 时间};

\node[
  anchor=center,
  circle,
  fill=black,
  inner sep=0,
  minimum size=4pt
] (tpoint)
  at ([xshift=-4.5cm]x1.west) {};

\node[anchor=east] (tlabel)
  at ([xshift=-0.1cm]tpoint)
  {\footnotesize $t$};

\draw[-,dotted]
  ([xshift=2pt]tpoint)
  --
  (xlabel.west);


% =========================================================
% CTMC 公式
% =========================================================

\node[anchor=south] (CTMC)
  at ([yshift=2pt]xbox.north)
  {
    \footnotesize
    CTMC:
    $
    \frac{d p_t(\mathbf{x})}{dt}
    =
    \sum_{\mathbf{y}\in V^n}
    p_t(\mathbf{y})
    {\color{blue}
      [\mathbf{Q}^{\mathrm{seq}}(t)]_{\mathbf{y}\mathbf{x}}
    }
    $
  };


% =========================================================
% 右上角：序列 x
% =========================================================

\node[anchor=center] (sx1)
  at ([yshift=4cm,xshift=2cm]xbox.north)
  {\footnotesize\texttt{A}};

\node[anchor=center] (sx2)
  at ([xshift=1.2*\stepsize]sx1)
  {\footnotesize\texttt{D}};

\node[anchor=center] (sx3)
  at ([xshift=1.2*\stepsize]sx2)
  {\footnotesize\texttt{B}};

\node[anchor=center] (sx4)
  at ([xshift=1.2*\stepsize]sx3)
  {\footnotesize\texttt{C}};

\node[anchor=center] (sx5)
  at ([xshift=1.2*\stepsize]sx4)
  {\footnotesize\texttt{C}};


% =========================================================
% 右上角：序列 y_3
% =========================================================

\node[anchor=center] (sy1)
  at ([yshift=-1.2cm]sx1.south)
  {\footnotesize\texttt{A}};

\node[anchor=center] (sy2)
  at ([xshift=1.2*\stepsize]sy1)
  {\footnotesize\texttt{D}};

\node[anchor=center] (sy3)
  at ([xshift=1.2*\stepsize]sy2)
  {\footnotesize\texttt{B}};

\node[anchor=center] (sy4)
  at ([xshift=1.2*\stepsize]sy3)
  {\footnotesize\texttt{D}};

\node[anchor=center] (sy5)
  at ([xshift=1.2*\stepsize]sy4)
  {\footnotesize\texttt{C}};


% =========================================================
% 右上角逐 token 转移箭头
% =========================================================

\draw[->,thin] (sy1.north) -- (sx1.south);
\draw[->,thin] (sy2.north) -- (sx2.south);
\draw[->,thin] (sy3.north) -- (sx3.south);
\draw[->,thin] (sy4.north) -- (sx4.south);
\draw[->,thin] (sy5.north) -- (sx5.south);


%  token 级生成矩阵元素

\node[anchor=north west] (qtokenlabel)
  at (sy4.south)
  {
    \scriptsize
    $Q^{\mathrm{tok}}_{\texttt{DC}}(t)$
  };


% =========================================================
% 修复部分：
% 蓝色外框和第四列 C/D 的浅橙色高亮
% =========================================================

\begin{pgfonlayer}{highlight}

% 右上角整个矩阵的蓝色外框
\node[
  draw=blue,
  rounded corners=1pt,
  line width=0.7pt,
  inner sep=3pt,
  fit=(sx1) (sy5) (qtokenlabel)
] (qbox) {};

% 第四列 C -> D 的颜色高亮
\node[
  fill=orange!20,
  rounded corners=0.6pt,
  inner xsep=2pt,
  inner ysep=1pt,
  fit=(sx4) (sy4)
] (qttokenbox) {};

\end{pgfonlayer}


% =========================================================
% 右上角蓝色弯曲箭头
% =========================================================

\draw[->,blue]
  ([yshift=-2pt]qbox.south)
  ..
  controls
    +(south:0.7cm)
    and
    +(north:0.7cm)
  ..
  ([xshift=-1.0cm,yshift=-1pt]CTMC.north east);


% =========================================================
% 右上角说明文字
% =========================================================

\node[
  anchor=north east,
  align=left
] (qlabel)
  at ([yshift=0]qbox.north west)
  {
    \footnotesize
    生成矩阵中 $(\mathbf{x},\mathbf{y}_3)$ 对应的
    \\[-2pt]
    \footnotesize
    矩阵元，记为
    $[\mathbf{Q}^{\mathrm{seq}}(t)]_{\mathbf{y}\mathbf{x}}$
    \\[-2pt]
    \footnotesize
    其中
    \\[-2pt]
    \footnotesize
    $\mathbf{x}=\texttt{ADBCC}$
    \\[-2pt]
    \footnotesize
    $\mathbf{y}_3=\texttt{ADBDC}$
  };

\end{scope}

\end{tikzpicture}

\end{document}